%======================================================================
%  BACHELOR'S THESIS PROPOSAL — the first kernel of the thesis
%  Ruhr-Universität Bochum (student) · MPI-SP (advisor)
%
%  This is the document sent to the supervisor. He asked for three
%  pages; it is laid out as a paper (ACM two-column, no ACM boilerplate)
%  and holds exactly the four parts agreed on 2026-09-04:
%      Background · Research Problem · Why It Matters · Methodology
%  plus the RQ and scope from his written input of 2026-09-05.
%
%  HARD CONSTRAINT: three full pages of content; the reference list
%  starts on page 4 and is not counted. After any edit, recompile and
%  check. If the text spills, cut content — never shrink font or margins.
%
%  The long, argued version (schedule, ten open points, full method)
%  lives in ../../99_Archive/_junk/02_Proposal/. Chapter 4 of the thesis
%  is written from it; this file is what grows into the thesis.
%
%  COMPILE:  latexmk   (pdflatex -> biber -> pdflatex -> pdflatex)
%======================================================================
\documentclass[sigconf,nonacm,natbib=false]{acmart}

%----------------------------------------------------------------------
%  Packages
%----------------------------------------------------------------------
\usepackage{csquotes}
\usepackage{booktabs}
\usepackage{array}
\usepackage{enumitem}
\usepackage{tikz}
\usetikzlibrary{arrows.meta}

\setlist{nosep,leftmargin=*}

% Ragged-right paragraph column for the narrow table cells; justified
% text in a 1.5 cm column opens wide gaps between words.
\newcolumntype{L}[1]{>{\raggedright\arraybackslash}p{#1}}

% Linux Libertine's small caps set "qu" as one ligature glyph that has no
% Unicode mapping, so the PDF text read "CONSEQENCE". Code names are set
% in small caps, so they must survive copy and search.
\DisableLigatures[q]{encoding=T1,shape=sc}

% Same shared bibliography as the thesis and the slides. acmart defaults
% to natbib + BibTeX; the repository uses biblatex + biber, hence natbib=false
% above (LaTeX reports it as an "unused option" — that is harmless).
\usepackage[backend=biber,style=numeric-comp,sorting=none,
            maxbibnames=2,minbibnames=1,maxcitenames=1,
            urldate=long]{biblatex}
\addbibresource{../../_shared/references.bib}
\renewcommand*{\bibfont}{\footnotesize}
\setlength{\bibitemsep}{0pt}

%======================================================================
\begin{document}

\title[Privacy Harm Communication in Health Tracking Apps]{Understanding
  Privacy Harm Communication\texorpdfstring{\\}{ }in Fitness and Health
  Tracking Apps}
\subtitle{Bachelor's Thesis Proposal}

\author{Hamza Sghaier}
\affiliation{%
  \institution{Ruhr University Bochum}
  \department{B.Sc.\ IT-Sicherheit / Informationstechnik}
  \city{Bochum}\country{Germany}}

\author{Kavous Salehzadeh Niksirat}
\authornote{Advisor.}
\affiliation{%
  \institution{Max Planck Institute for Security and Privacy (MPI-SP)}
  \city{Bochum}\country{Germany}}

\renewcommand{\shortauthors}{Sghaier and Salehzadeh Niksirat}

%---------------------------- Abstract --------------------------------
\begin{abstract}
Fitness and period trackers process data that can reveal a person's health
and sex life, both special categories under the GDPR, yet what their
providers \emph{say} about the resulting risks is barely studied. Prior work
asks whether providers disclose their data practices truthfully and what
users believe; to our knowledge, it does not ask which privacy \emph{harms}
providers name, minimise, or leave out. We propose a content analysis of the
privacy communication of 50 highly rated fitness and fertility apps on the
German Google Play and Apple App Store --- privacy policies, store listings,
provider marketing, and the app interface itself --- coded against harm trees
that trace data flows to harms, with LLM assistance that is validated for what it \emph{misses} as
well as what it finds.
\end{abstract}

\keywords{privacy harms, health trackers, privacy communication, privacy-enhancing technologies, content analysis}

\maketitle

%---------------------------- 1. Background ---------------------------
\section{Background}

In Germany, 73\,\% of smartphone users use health apps for fitness, nutrition,
or general wellbeing~\parencite{bitkom2025health}. Fitness, sleep, and menstrual
trackers accumulate three kinds of data at once: attributes users enter
(weight, symptoms, mood, sexual activity), continuous sensor streams (steps,
heart rate, sleep, location), and system context (identifiers, timestamps, usage
patterns). Under Article~9(1) GDPR, data concerning health and data concerning
a person's sex life are special categories of personal data. Cycle records
concern health and often sex life as well; activity data such as step counts
qualifies once it reveals information about health status (Article~4(15),
Recital~35)~\parencite{gdpr2016}.

A privacy \emph{violation} is not the same as a privacy \emph{harm}. This
thesis takes its harm model from Knabenhans et al., who analyse the European
Digital Identity Framework with \emph{harm trees}~\parencite{knabenhans2026harms}.
A tree has three levels. At the bottom are \emph{system components}: information
the system leaks, and requirements it imposes on users. Components combine into
adversarial \emph{capabilities}: distinguishing users, tracking or profiling
them, and inferring their attributes. Capabilities lead to \emph{harms}, sorted
into seven types --- physical, economic, psychological, autonomy,
discrimination, reputational, and societal --- that extend Citron and Solove's
typology~\parencite{citron2021harms}. Each component is traced to its origin:
the core functionality, the digital context, or an implementation choice.
Privacy-enhancing technologies (PETs) can remove only the last kind. For every
harm Knabenhans et al.\ study, a pathway remains that no PET can close, and
presenting PETs as the answer risks a false sense of safety.

The model transfers directly to trackers. Inference there is not hypothetical:
a survey of 236 studies reports that tracker data can be mined for private
attributes, personality traits among them~\parencite{niksirat2024wearable}. In
a cycle tracker, inference is the core functionality itself: predicting the
fertile window is what the app is for, so no PET can remove that capability.
Martin et al.\ examine three principles behind big-tech PET deployments ---
limiting third-party access, using inferences instead of raw data, keeping data
on the device --- and argue that each can be satisfied while the
privacy-violating behaviour continues~\parencite{martin2025cookies}. Whether
tracker providers invoke the same principles is what RQ3 asks.

Providers communicate about privacy on four surfaces: the privacy policy, the
store listing, their marketing (website, blog, advertising), and the app
interface. All four are regulated --- misleading claims in listings and
advertising, for instance, fall under unfair-commercial-practices
law~\parencite{ucpd2005} --- but only one has prescribed content: the notice
required by Articles 13--14 GDPR, conventionally the privacy policy. It must
name the controller, purposes, legal basis, recipients, transfers, retention
period, and the user's rights. It describes \emph{data flows}, not
\emph{consequences}.

Inference and exposure are covered only in part. Inference is processing, so a
provider that profiles its users must name profiling as a purpose. It need not
say what it infers: \enquote{to personalise your experience} is enough, even if
it estimates pregnancy or illness. Users learn what was inferred only on
request (Article~15). For exposure, the provider must assess the risks
(Articles~32 and~35) but need not share the assessment. It must state the
\enquote{likely consequences} of a breach only after a high-risk breach
(Article~34(2)). Beyond that, consequences must be stated only for refusing to
provide data and for solely automated decisions (Articles 13(2)(e),
13(2)(f))~\parencite{gdpr2016}. Recitals~39 and~60 and the EDPB-endorsed
guidelines ask for more~\parencite{wp292018transparency}, but neither is
binding. \textbf{Outside solely automated decisions, no binding provision
requires a provider to tell users in advance what could be inferred from their
data, or what would follow if it were exposed.} What providers say about harm
is voluntary; this thesis examines what they say.

%---------------------------- 2. Research Problem ---------------------
\section{Research Problem}

A large literature asks whether providers tell the truth. Huckvale et al.\ found
29 of 36 depression and smoking-cessation apps sending data to Google or
Facebook services while only 12 disclosed it~\parencite{huckvale2019assessment};
a dynamic study of nearly 5{,}000 Android apps found 81\,\% misrepresenting
their practices in the Data Safety section~\parencite{arkalakis2024abandon};
Apple's privacy nutrition labels have been analysed at
scale~\parencite{li2022ios}; menstrual and fertility trackers have been
examined through policy analysis, usability inspection, and network
measurement~\parencite{shipp2020period,malki2024exploring,hudig2025intimate}.
A second literature asks what users believe: fitness tracker users know
\emph{what} is collected but not \emph{how} it is used, and rate harmful
scenarios as plausible yet unlikely~\parencite{gabriele2020understanding}; US
period tracker users report high concern together with feeling uninformed and
powerless~\parencite{cao2024deleted}. The closest precedent for coding stance,
an analysis of 161 breach notifications, shows companies hedging what they say
\emph{after} a breach~\parencite{zou2019might}.

To our knowledge, none maps which \emph{harms} providers name
\emph{before} anything has gone wrong. That is the gap: \textbf{we lack a
systematic account of what harms service providers communicate} --- which
types are acknowledged, minimised, or ignored, and how PETs are framed (secrecy
versus harm prevention). In harm-tree terms, the GDPR's information duties and
the store disclosure schemas ask only for the bottom level: data flows. We ask
whether provider text ever climbs higher --- to a \textbf{capability} such as
inference, or to a \textbf{harm} --- and what it says about two things outside
those duties: \textbf{breach consequences} and \textbf{incident handling}.

\noindent\textbf{RQ.} \emph{What types of downstream privacy harms do
health tracker service providers explicitly acknowledge or downplay when
describing their privacy practices or PETs, and how do they frame the
protection these technologies offer?}

\begin{description}[leftmargin=1.5em,itemsep=1pt,topsep=2pt,labelindent=0pt,
                    font=\normalfont\bfseries]
  \item[RQ1] \textit{Coverage.} Which harms, capabilities, and components of
    the domain's harm tree are named, which are absent, and how often does a
    text reach the harm level at all?
  \item[RQ2] \textit{Stance.} Is a mentioned harm acknowledged, hedged,
    minimised, or deflected onto the user?
  \item[RQ3] \textit{PET framing.} Which protections are claimed, which
    pathways can they close given their origin, and are they presented as
    preventing harms whose pathways they cannot close?
  \item[RQ4] \textit{Consistency.} Does the answer change between policy, store
    listing, marketing, and interface?
\end{description}

\noindent\textbf{Scope.} \enquote{Health trackers} here means fitness and
activity tracking (sleep, heart rate, nutrition included) and fertility/period
tracking, on the German storefronts of both stores, analysed through the app
--- on smartphones, and for wearables such as smart watches through their
companion app. Out of scope: a user study (a possible later phase), verifying
whether claims are \emph{true}, grading legal compliance, and DiGA (reimbursable
digital health applications).

%---------------------------- 3. Why it matters -----------------------
\section{Why the Problem Is Important}

\textbf{Users.} Installing a cycle tracker means weighing a tangible benefit
against a risk that has to be imagined. If the communication never names a
consequence, users must supply it themselves, and evidence suggests they
discount it: fitness-tracker users rate harmful scenarios as
unlikely~\parencite{gabriele2020understanding}. For fertility data the stakes
are concrete: after Roe v.~Wade was overturned, some US period-tracker users
deleted their apps for fear their data could be used against
them~\parencite{cao2024deleted}. A user told that an app does not share raw data
may conclude they are protected; if the provider derives and acts on an
inference instead, that conclusion is wrong and nothing in the text said so.
\textbf{Providers and designers.} A map of which harms are named, where, and
in what terms shows what their communication omits.
\textbf{Regulators.} If text rarely climbs above data flows across a
segment handling Article~9 data, that is a concrete finding about what a
transparency regime built on data-practice disclosure leaves out.
\textbf{Research.} Harm typologies were built for courts and for system
analysis; to our knowledge, neither has been applied to the text providers put
in front of users.

%---------------------------- 4. Methodology --------------------------
\section{Methodology}

A structured content analysis of privacy communication, plus a manual
inspection of the interfaces that carry part of it (Figure~\ref{fig:pipeline}).
Every code is anchored to a verbatim quote with source, location, and capture
date; every instrument --- sampling protocol, codebook, each LLM prompt --- is
versioned, and every model run is logged with prompt version, model identifier,
date, and raw output, so any figure can be traced to a dated source.

\textbf{Literature review.} A logged-search review tests the claimed gap rather
than assuming it, in three areas: \emph{privacy harms}
\parencite{solove2006taxonomy,calo2011boundaries}, \emph{health tracker
privacy}, and \emph{methods}, from policy
annotation~\parencite{wilson2016creation,zimmeck2019maps} to content
analysis~\parencite{krippendorff2019content}.

\begin{figure*}[t]
\centering
\begin{tikzpicture}[
  font=\scriptsize,
  bx/.style ={draw, rounded corners=2pt, align=center, text width=29mm,
              minimum height=13mm, inner sep=2.5pt},
  arr/.style={-{Stealth[length=1.6mm]}, semithick}
]
  \hyphenpenalty=10000 \exhyphenpenalty=10000
  \node[bx] (a) at (-7.2,0) {\textbf{Sample}\\50 apps, German stores,\\stratified 30/20};
  \node[bx] (b) at (-3.6,0) {\textbf{Corpus}\\policies, listings, marketing,\\12 walkthroughs};
  \node[bx] (c) at ( 0.0,0) {\textbf{Coding}\\against per-domain\\harm trees};
  \node[bx] (d) at ( 3.6,0) {\textbf{Validation}\\precision, $\kappa$ with prevalence,\\recall with interval};
  \node[bx] (e) at ( 7.2,0) {\textbf{Answers}\\RQ1--RQ4, each code\\backed by a quote};
  \draw[arr] (a) -- (b); \draw[arr] (b) -- (c);
  \draw[arr] (c) -- (d); \draw[arr] (d) -- (e);
\end{tikzpicture}
\caption{The pipeline, from app sample to answers.}
\Description{Five boxes in a row joined by arrows: sample, corpus, coding,
validation, answers.}
\label{fig:pipeline}
\end{figure*}

\textbf{Sampling.} Google Play and the Apple App Store, German storefronts. A
frozen keyword list covers both domains in German, plus English terms used
untranslated in German listings (e.g.\ \emph{Schrittzähler, Workout; Periode,
Eisprung, period tracker}).
\enquote{Highest rated} is fixed in advance: the first 20 results per keyword
are recorded, the pooled set filtered to apps with $\geq 1{,}000$ ratings (a
heuristic floor, so that averages rest on enough votes), then ranked by average
rating, rating count as tie-break; an app on both stores counts once.
Target: \textbf{50 apps}, stratified 30 fitness / 20 fertility, reported per
stratum. A five-app pilot tests the instruments; its apps stay in the sample
but are kept out of the validation set. Included apps collect health or body
data, are free to install, and have a reachable policy; DiGA, device companions
with no tracking of their own, and duplicates are excluded. An app with no
reachable policy is recorded as a finding. The ranked lists are stored with the
capture date, so the sample can be reconstructed.

\textbf{Material.} Per app: privacy policy, store description and structured
disclosure, and the provider's marketing --- website, privacy blog posts, and
advertising. German versions are coded, English ones checked against them.
Each document is archived with URL, capture date, hash, and app version, then
split into coding units at paragraph boundaries. Prevalence is reported per
app, so that a 40-page policy does not outweigh a three-line blurb. Absent
advertising is not a finding, since ad delivery cannot be reproduced.

\textbf{Walkthrough.} Twelve apps --- per domain, the three highest-ranked
eligible apps on Android and on iOS, excluding apps that cannot be set up
without a paired device --- are inspected manually with the cognitive
walkthrough~\parencite{wharton1994cognitive}, adapted so that the tasks are
\emph{privacy} tasks, from a fresh install and screen-recorded. The tasks: find
what is collected; find what is \emph{derived}; find what happens after a
breach; reach the policy; stop third-party sharing; export; delete. Each step
is recorded as success or failure with the reason, and the privacy text on its
screens is coded with the same codebook.

\textbf{Reference harm trees.} Before coding, one harm tree is built per
domain, following Knabenhans et al.'s procedure: harms are taken from the
tracker literature~\parencite{niksirat2024wearable,shipp2020period,cao2024deleted}
and the seven types, and each pathway's components are labelled by origin. In a
cycle tracker, for instance, logged cycles and the predicted fertile window are
\emph{core}; device identifiers and platform metadata are \emph{context};
advertising SDKs, cloud sync, and retention periods are \emph{implementation}.
The trees fix in advance what a complete account of harm could contain, so an
absence is measured against a stated, versioned baseline.

\textbf{Coding.} Each unit is mapped onto the tree on three axes
(Table~\ref{tab:axes}); stance is coded only where a harm is named. Every category gets a definition, a positive example, and a
mandatory counter-example; without the counter-example a category has no
boundary, and human and model alike drift. Text is coded as written, never as
what the provider is assumed to mean; a unit may carry several codes; when
unsure, a unit is coded and flagged, never dropped; any change to a category
produces a new codebook version with date and reason. The same codebook binds
the model and the human coder; disagreements are adjudicated against the
verbatim span and logged.

\begin{table}[h]
\caption{The three coding axes.}
\label{tab:axes}
\small
\begin{tabular}{@{}L{1.55cm}L{6.3cm}@{}}
\toprule
\textbf{Axis} & \textbf{Values} \\
\midrule
A\enspace Pathway &
  highest tree level reached (component, capability, harm); capability
  (distinguishing, tracking/profiling, inference); harm type (seven);
  probes \textsc{breach-consequence}, \textsc{incident-handling} \\
B\enspace Stance &
  acknowledged, hedged, minimised, deflected; at app level also absent \\
C\enspace Protection &
  measure claimed (incl.\ Martin et al.'s principles); origin of the
  component it removes; data-flow vs.\ harm-prevention framing;
  \textsc{overclaim}; reassurance with no verifiable commitment \\
\bottomrule
\end{tabular}
\end{table}

\textsc{overclaim} marks a measure presented as preventing a harm whose
reference tree keeps a core or context pathway --- the false sense of safety
Knabenhans et al.\ warn about. Stance is not part of their model; it is adapted
from the hedging found in breach notices~\parencite{zou2019might}.
Table~\ref{tab:example} codes a constructed sentence: a text can be accurate
and still silent on harm.

\begin{table}[h]
\caption{A constructed example of a coded unit.}
\label{tab:example}
\small
\begin{tabular}{@{}L{0.9cm}L{6.9cm}@{}}
\toprule
\textbf{Unit} & \enquote{We never share your raw data with third parties;
  all analysis happens on your device.} \\
\midrule
A & component level only; no capability or harm named \\
B & not coded, since no harm is named \\
C & third-party limitation and on-device processing; both remove
  implementation components; framed as data flow. In a cycle tracker the
  inference pathway is core, so a claim that this \emph{prevents} harm
  would be coded \textsc{overclaim} \\
\bottomrule
\end{tabular}
\end{table}

\textbf{LLM assistance and validation.} The corpus is too large to code by
hand, so coding is LLM-assisted: one versioned prompt per axis, a pinned model
(chosen on the pilot, its identifier frozen before the main run) at temperature
0, and structured output in which every code carries a verbatim span --- a code
without one is discarded. Temperature 0 does not guarantee identical output from
a hosted model, so each unit is run three times and a code is kept when at
least two runs return it. Existing tools~\parencite{harkous2018polisis,rodriguez2024llm}
code \emph{data practices}, not harms, so the pipeline is purpose-built.
\textbf{Calibration is deliberately asymmetric.} The headline claims are claims of
\emph{absence}, and the two error types are not equally dangerous: a false
positive inflates a count that ships with its quote and can be checked, whereas
a false negative manufactures the very silence the thesis sets out to establish.
The author therefore codes a validation set by hand, blind to the model's
output. On a random 20\,\% of units, per-code precision and model--human
Cohen's $\kappa$~\parencite{cohen1960coefficient} are reported, $\kappa$
\emph{with} prevalence, since it is depressed when a code is
rare~\parencite{feinstein1990high}. A random sample holds few instances of rare
codes, so for the \textbf{recall audit} the set is enriched: a lexical sweep,
independent of the model, flags harm-adjacent vocabulary in both languages, and
flagged units are added until each harm code has 30 hand-confirmed instances
or none remain; every uncoded hit in the whole corpus is read by hand. An absence is
thus never reported as \enquote{the provider does not mention X} but as
\enquote{no mention was found by an instrument whose measured recall for X is
$r$}, with a 95\,\% interval, or as unestimable where too few instances exist.

\textbf{Analysis.} Descriptive and per domain; only the validation estimates
carry intervals (Table~\ref{tab:rq}). The headline for RQ4 is how often a code
present in the policy is absent from the screen where the user decides.

\begin{table}[h]
\caption{From research question to result.}
\label{tab:rq}
\small
\begin{tabular}{@{}L{0.75cm}L{3.2cm}L{3.7cm}@{}}
\toprule
\textbf{RQ} & \textbf{Evidence} & \textbf{Result} \\
\midrule
RQ1 & all units, all surfaces & tree nodes reached, per app \\
RQ2 & units that name a harm & stance distribution, as counts \\
RQ3 & all surfaces & measures by origin; overclaims \\
RQ4 & all; interface for 12 apps & pairwise comparison of code sets \\
\bottomrule
\end{tabular}
\end{table}

\textbf{Ethics.} No human subjects; walkthroughs use synthetic data and study
accounts; store terms are checked before any automated retrieval.

\textbf{Limitations.} \emph{Absence is not intent:} the study establishes that a
harm is not mentioned, never \emph{why}, and attributes no motive.
\emph{Snapshot:} every claim is indexed to the capture date.
\emph{Selection:} the German storefront, rating floor, and search ranking
favour large providers; the findings describe what users most likely meet.
\emph{Trees:} the reference trees are the author's construction, so coverage is
measured against them, not against all possible harms. \emph{One coder:} $\kappa$ compares
model and author, so no inter-coder reliability is
claimed~\parencite{mcdonald2019reliability}. \emph{Recall:} $r$ may overstate
recall for implicit mentions the sweep did not anticipate.

%---------------------------- References ------------------------------
\clearpage
% acmart balances the last page only at \end{document}, when the second
% column has already started; calling \balance here evens out the
% reference page.
\balance
\printbibliography[title={References}]

\end{document}
